% Paste this block into Chapter 1, directly below
%   \section{Proposed Solution}\label{sec:solution}
%   % high-level representative fig 1.1
% Files: copy out/digital_twin_light.png to figures/digital_twin_light.png (use digital_twin_dark.png for slides) and
% add out/references_additions.bib to your bibliography (\bibliography{...}).
% Requires: \usepackage{graphicx}, \usepackage{textcomp} or any package providing \dag, and \usepackage{url}/hyperref.

\begin{figure}[t]
  \centering
  \includegraphics[width=\textwidth]{figures/digital_twin_light.png}
  \caption{Digital twin of the proposed AI-based DOFS perimeter security system for Makkah during Hajj.
  The fibre follows the boundary of the Haram and Arafat and is coloured by vibration activity; bright segments
  are light pulses launched by each interrogator unit (IU) along the fibre segment it serves. Dashed lines mark the
  10\,km section covered by the simulated trace on the right. A walking intruder (red) and a vehicle on a road (amber)
  are both detected because their activity exceeds the detection threshold. In the proposed system, an AI classifier
  then decides: the intruder is classified as human and raises an alarm that sends the nearest UAV and patrol, whereas
  the vehicle is classified as a vehicle on a road and is only logged (marked \dag{}). This classification step is
  part of the proposed system and is not simulated. Shaded areas mark weak coverage, computed from terrain slope,
  directional CCTV visibility and distance to roads. The inset enlarges the rectangle drawn on the map around the
  Mashaer (Jamarat, Mina, Muzdalifah, Arafat); north is up in both views.
  Fibre length and interrogator count are computed from the boundary geometry and the specification of the reference
  interrogator~\cite{eoss_das_datasheet} (70\,km distance range per channel, one built-in channel):
  $N_\mathrm{IU}=\lceil 145.5/(70\times 1)\rceil=3$. CCTV, UAV and patrol numbers and positions, the C2 location and the road
  entry points used as checkpoint positions (marked \dag{}) are proposed placements, not real sites. UAV altitude and
  cone size are not to scale, the vertical exaggeration is $2\times$, and the event hot spots on the fibre are
  widened for visibility.
  The trace is the output of the $\varphi$-OTDR simulation~\cite{das_simulation} for a 10\,km section with default
  parameters, not measured data. Its walking event is the intruder on the map and its vehicle event is the vehicle;
  the simulated digging event and calibration sine are not shown on the map. The detection threshold follows the
  simulator's rule (median plus $8\times1.4826\times$ MAD of the activity).
  Hajj 2026 pilgrim count: GASTAT~\cite{gastat2026hajjnews}.
  Boundary and roads: OpenStreetMap contributors~\cite{osm_contributors}. Terrain: produced using Copernicus WorldDEM-30
  \copyright{} DLR e.V. 2010--2014 and \copyright{} Airbus Defence and Space GmbH 2014--2018 provided under COPERNICUS by the
  European Union and ESA; all rights reserved~\cite{copernicus_dem_glo30}.
  Map image is the intellectual property of Esri and is used herein under license.
  Copyright \copyright{} 2026 Esri and its licensors. All rights reserved.~\cite{esri_world_imagery}}
  \label{fig:digital-twin}
\end{figure}
